\subsection{Unendlich viele Lösungen parametrisieren}
\label{subsec:linear_systems_parametric_solution_sets}

Besitzt ein lösbares LGS freie Variablen, werden diese als Parameter gewählt.
Die Pivotvariablen ergeben sich dann durch Rückwärtseinsetzen.

\begin{ExampleBox}{Lösungsmenge mit einem Parameter}{linear_systems_parametric_solution}
  Nach der Vorwärtselimination sei
  \[
    \left[
      \begin{array}{rrr|r}
        2&1&-1& 3\\
        0&7&-5&-7\\
        0&0& 0& 0
      \end{array}
    \right].
  \]
  Es gilt
  \[
    \operatorname{rang}(A)=\operatorname{rang}([A\mid b])=2<3.
  \]
  Somit ist $x_3$ frei. Die Wahl $x_3=7t$ vermeidet Brüche. Dann folgt
  \begin{align*}
    7x_2-5x_3&=-7
      &\Longrightarrow\quad x_2&=-1+5t,\\
    2x_1+x_2-x_3&=3
      &\Longrightarrow\quad x_1&=2+t.
  \end{align*}
  Daher ist
  \[
    L=
    \left\{
      \begin{pmatrix}2\\-1\\0\end{pmatrix}
      +t\begin{pmatrix}1\\5\\7\end{pmatrix}
      \;\middle|\;t\in\R
    \right\}.
  \]
  Der erste Vektor ist eine spezielle Lösung; der zweite gibt die Richtung der
  Lösungsgeraden an.
\end{ExampleBox}

\begin{KeyPointBox}{Zahl der freien Variablen}{linear_systems_number_free_variables}
  Ist ein LGS lösbar, so gilt
  \[
    \boxed{\text{Anzahl freier Variablen}
      =n-\operatorname{rang}(A)}.
  \]
  Jede freie Variable liefert einen unabhängigen Parameter der Lösungsmenge.
\end{KeyPointBox}

Für ein homogenes LGS $Ax=0$ folgt unmittelbar:
\[
  \operatorname{rang}(A)=n
    \Longleftrightarrow \text{nur die triviale Lösung},
  \qquad
  \operatorname{rang}(A)<n
    \Longleftrightarrow \text{unendlich viele Lösungen}.
\]
